\begin{center}
    \textbf{BAB I}\\
    \textbf{PENDAHULUAN}
    \addcontentsline{toc}{chapter}{BAB I. PENDAHULUAN}
\end{center}

\textbf{A. Latar Belakang}
\addcontentsline{toc}{section}{A. Latar Belakang}

\hspace*{1cm}Pendidikan merupakan salah satu faktor utama dalam pembangunan suatu bangsa. Dalam upaya meningkatkan mutu pendidikan, diperlukan inovasi dalam metode pembelajaran yang mampu memenuhi tantangan zaman serta kebutuhan peserta didik. Salah satu tantangan yang sering dihadapi dalam pembelajaran matematika adalah kesulitan peserta didik dalam memahami konsep-konsep abstrak, khususnya pada materi bangun ruang sisi datar. Materi ini sulit dipahami karena umumnya hanya disampaikan melalui media dua dimensi, seperti buku teks atau gambar di papan tulis, sehingga peserta didik kesulitan dalam memvisualisasikan bentuk tiga dimensi secara nyata \citep{battista1999}.

\hspace*{1cm}Kemajuan teknologi di era digital memberikan peluang besar untuk menghadirkan inovasi dalam dunia pendidikan. Salah satu teknologi yang memiliki potensi besar adalah Augmented Reality (AR) \citep{saidin2015}. Teknologi AR memungkinkan penggabungan objek virtual dengan dunia nyata secara interaktif dan realistis. Dalam pembelajaran matematika, penggunaan AR diharapkan dapat meningkatkan kemampuan visualisasi, pemahaman, dan pengalaman belajar peserta didik terhadap materi bangun ruang \citep{syahputra2024}.

\hspace*{1cm}Di SMP Negeri 3 Kalikajar, proses pembelajaran matematika masih didominasi oleh metode konvensional. Guru menggunakan buku teks dan penjelasan di papan tulis sebagai media utama. Namun, hasil evaluasi pembelajaran menunjukkan bahwa banyak peserta didik masih mengalami kesulitan dalam memahami materi bangun ruang sisi datar. Kondisi ini menunjukkan perlunya pengembangan media pembelajaran yang lebih inovatif dan interaktif guna membantu peserta didik memahami konsep secara lebih mendalam.

\hspace*{1cm}Penggunaan media pembelajaran berbasis Augmented Reality (AR) diharapkan dapat menjadi solusi yang efektif dalam mengatasi kesulitan pembelajaran matematika. AR sebagai media pembelajaran interaktif diharapkan dapat meningkatkan pemahaman peserta didik terhadap konsep-konsep matematika, khususnya pada materi bangun ruang sisi datar \citep{nincarean2013mar,mustaqim2016pemanfaatan,qumillaila2017eksresi}. Media pembelajaran berbasis AR mengintegrasikan elemen digital dengan dunia nyata melalui visualisasi 3D, animasi, dan interaksi langsung yang meningkatkan keterlibatan peserta didik dalam pembelajaran. Peserta didik menunjukkan antusiasme dan partisipasi aktif yang lebih tinggi dibandingkan dengan metode konvensional \citep{arinda2023}. Dengan memanfaatkan teknologi ini, diharapkan pemahaman peserta didik terhadap materi bangun ruang sisi datar dapat meningkat secara signifikan, sehingga hasil belajar menjadi lebih optimal.

\hspace*{1cm}Berdasarkan uraian tersebut, penelitian ini bertujuan untuk mengembangkan media pembelajaran berbasis Augmented Reality pada materi bangun ruang sisi datar di SMP Negeri 3 Kalikajar. Penelitian ini diharapkan dapat memberikan kontribusi nyata dalam meningkatkan kualitas pembelajaran matematika, serta menjadi referensi bagi pengembangan media pembelajaran berbasis teknologi di masa yang akan datang.\\


\textbf{B. Identifikasi Masalah}
\addcontentsline{toc}{section}{B. Identifikasi Masalah}

\hspace*{1cm}Berdasarkan uraian pada latar belakang, dapat diidentifikasi beberapa permasalahan sebagai berikut:
\begin{enumerate}
    \item Peserta didik mengalami kesulitan dalam mempelajari materi geometri, khususnya bangun ruang sisi datar, karena mereka kesulitan memvisualisasikan objek geometri tiga dimensi.
    \item Masih banyak peserta didik kelas VIII SMP Negeri 3 Kalikajar yang memperoleh nilai ulangan harian pada materi bangun ruang sisi datar di bawah Kriteria Ketuntasan Minimal (KKM).
    \item Pembelajaran geometri bangun ruang sisi datar di SMP Negeri 3 Kalikajar masih didominasi oleh metode tradisional, seperti penyampaian materi melalui gambar di papan tulis.
\end{enumerate}

\textbf{C. Pembatasan Masalah}
\addcontentsline{toc}{section}{C. Pembatasan Masalah}

\hspace*{1cm}Berdasarkan identifikasi masalah, penelitian ini difokuskan pada pengembangan media pembelajaran berbasis Augmented Reality (AR) pada materi bangun ruang sisi datar guna mempermudah peserta didik dalam memvisualisasikan objek geometri tiga dimensi. Penelitian ini bertujuan untuk menguji tingkat kelayakan media yang dikembangkan, tanpa mengukur pengaruh media tersebut terhadap hasil belajar peserta didik.\\

\textbf{D. Rumusan Masalah}
\addcontentsline{toc}{section}{D. Rumusan Masalah}

\hspace*{1cm}Berdasarkan latar belakang dan identifikasi masalah, maka rumusan masalah dalam penelitian ini adalah sebagai berikut:
\begin{enumerate}
    \item Bagaimana prosedur pengembangan media pembelajaran berbasis Augmented Reality pada materi bangun ruang sisi datar?
    \item Bagaimana tingkat kelayakan media pembelajaran berbasis Augmented Reality pada materi bangun ruang sisi datar?
\end{enumerate}

\textbf{E. Tujuan Pengembangan}
\addcontentsline{toc}{section}{E. Tujuan Pengembangan}

\hspace*{1cm}Berdasarkan rumusan masalah di atas, maka tujuan dari penelitian ini adalah sebagai berikut:
\begin{enumerate}
    \item Mengetahui prosedur pengembangan media pembelajaran berbasis Augmented Reality pada materi bangun ruang sisi datar.
    \item Mengetahui tingkat kelayakan media pembelajaran berbasis Augmented Reality pada materi bangun ruang sisi datar.
\end{enumerate}

\textbf{F. Spesifikasi Produk yang Dikembangkan} \\
\addcontentsline{toc}{section}{F. Spesifikasi Produk yang Dikembangkan}
\hspace*{1cm}Produk yang dikembangkan oleh peneliti memiliki spesifikasi sebagai berikut:
\begin{enumerate}
    \item Produk berupa website yang dapat diakses melalui ponsel maupun laptop secara online.
    \item Produk mengintegrasikan materi bangun ruang sisi datar dengan teknologi Augmented Reality.
    \item Produk menampilkan objek AR melalui pemindaian kode QR (barcode).
\end{enumerate}

\textbf{G. Manfaat Pengembangan}
\addcontentsline{toc}{section}{G. Manfaat Pengembangan}

\hspace*{1cm}Hasil dari penelitian ini diharapkan dapat memberikan manfaat dalam dunia pendidikan, baik secara teoritis maupun praktis, sebagai berikut:

\begin{enumerate}
    \item \textbf{Manfaat Teoretis} \\
    Penelitian ini diharapkan memberikan kontribusi pengetahuan dalam bidang pendidikan matematika, khususnya terkait pengembangan media pembelajaran berbasis Augmented Reality.

    \item \textbf{Manfaat Praktis}
    \begin{enumerate}
        \item \textbf{Bagi Peserta Didik} \\
        Media ini diharapkan dapat membantu peserta didik dalam memvisualisasikan objek tiga dimensi, memahami materi bangun ruang sisi datar dengan lebih baik, serta meningkatkan minat belajar matematika.
        
        \item \textbf{Bagi Pendidik} \\
        Media ini dapat menjadi alternatif dalam menyampaikan materi bangun ruang sisi datar dan motivasi bagi pendidik untuk mengembangkan media pembelajaran berbasis Augmented Reality pada materi lainnya.
    \end{enumerate}
\end{enumerate}

\textbf{H. Asumsi dan Keterbatasan Pengembangan}
\addcontentsline{toc}{section}{H. Asumsi dan Keterbatasan Pengembangan}

\begin{enumerate}
    \item \textbf{Asumsi Pengembangan} \\
    Pengembangan media pembelajaran berbasis Augmented Reality ini didasarkan pada beberapa asumsi sebagai berikut:
    \begin{enumerate}
        \item Media pembelajaran berbasis Augmented Reality ini diharapkan dapat digunakan dalam pembelajaran matematika, khususnya pada materi bangun ruang sisi datar.
        \item Peserta didik akan lebih tertarik untuk belajar matematika ketika menggunakan media pembelajaran berbasis Augmented Reality.
        \item Penyajian bangun ruang sisi datar dalam bentuk virtual melalui teknologi AR dapat membantu peserta didik dalam memvisualisasikan konsep secara lebih nyata.
    \end{enumerate}

    \item \textbf{Keterbatasan Pengembangan} \\
    Penelitian pengembangan media pembelajaran berbasis Augmented Reality ini memiliki beberapa keterbatasan, antara lain:
    \begin{enumerate}
        \item Materi yang dikembangkan terbatas pada materi bangun ruang sisi datar.
        \item Subjek penelitian terbatas pada peserta didik kelas VIII SMP Negeri 3 Kalikajar.
        \item Media pembelajaran disusun berdasarkan Kurikulum 2013.
        \item Produk hanya dapat diakses secara online.
    \end{enumerate}
\end{enumerate}